% Chapter 2

\chapter{Reentrancy and Security Discussion}
\label{Chapter2}
\setlength{\parindent}{0pt}

%----------------------------------------------------------------------------------------
\section{Gas Evaluation}

Gas consumption was evaluated using the Hardhat testing environment. The analysis focused on the main protocol operations, including deposits, proposal approval, loan creation, repayments, and compensation claims.

To improve efficiency, the contracts minimize storage writes and use \texttt{immutable} variables for loan parameters that never change after deployment. The modular architecture, where each approved proposal deploys a dedicated \texttt{Loan} contract, increases deployment cost but simplifies the main contract and improves maintainability.

Some functions iterate over the contributor list to compute voting weights, available liquidity, and fund allocation, causing gas costs to grow with the number of contributors. This trade-off was considered acceptable for the project. Oracle integration is also gas-efficient because Bitcoin data are processed off-chain and only the required balances are stored on Ethereum.

%----------------------------------------------------------------------------------------
\section{Reentrancy Protection}

The project includes both a secure implementation (\texttt{LendingService.sol}) and an intentionally vulnerable version (\texttt{LendingServiceVulnerable.sol}) used to demonstrate a reentrancy attack.

The vulnerability affects the \texttt{claim\_compensation()} function, where the vulnerable contract transfers Ether before updating its internal state. A malicious contract can exploit this behavior by re-entering the function from its \texttt{receive()} callback while the previous execution is still in progress, allowing multiple compensation claims before the contributor's accounting information is updated.

The secure implementation follows the Checks--Effects--Interactions pattern by updating the contract state before performing external transfers, preventing recursive calls from reusing outdated accounting information.

The attack is implemented in \texttt{ReentrancyAttacker.sol} and validated through dedicated Hardhat tests, which confirm that the vulnerable contract can be exploited while the protected implementation resists the same attack.

%----------------------------------------------------------------------------------------
\section{Discussion of Malicious Strategies}

Besides reentrancy, other potential attack scenarios were considered during the design of the protocol.

The voting mechanism is based on contributors' available liquidity, making voting influence proportional to the amount of funds supplied rather than assigning equal voting power to each address. This aligns voting power with the economic stake committed to the protocol.

Loan approval also depends on multiple independent conditions. In addition to a positive vote, sufficient liquidity must be available and the Bitcoin collateral value provided by the oracle must satisfy the required collateralization threshold. Combining these conditions prevents proposals from being approved solely through voting.

The protocol relies on a trusted oracle owner to submit correct Bitcoin balance updates. Since blockchain data are obtained off-chain, the oracle remains a trusted component whose correctness directly affects collateral evaluation.

Although these measures improve the robustness of the system, they do not eliminate every possible attack vector. As with most decentralized applications, the overall security also depends on the correct behaviour of external components and administrative accounts.
